\chapter*{Introducción}
\addcontentsline{toc}{chapter}{Introducción}

El crecimiento sostenido de la infraestructura digital ha convertido el consumo energético del software en una preocupación central para la ingeniería de sistemas contemporánea. Durante años, las decisiones de diseño arquitectónico se evaluaron principalmente en función del rendimiento, la escalabilidad y la mantenibilidad, dejando en segundo plano su impacto sobre el consumo de recursos computacionales y, por extensión, sobre el consumo energético asociado a su ejecución \cite{kern2018sustainable}. Sin embargo, el incremento del volumen de datos, la proliferación de servicios en la nube y la adopción masiva de arquitecturas distribuidas han hecho evidente que la forma en que se organiza un sistema de software --monolítica o modular-- puede tener consecuencias medibles sobre la cantidad de recursos de cómputo que dicho sistema demanda \cite{araujo2024energy}.

La arquitectura monolítica, caracterizada por la concentración de la lógica de negocio en una única unidad desplegable, ha sido históricamente la opción predominante en el desarrollo de aplicaciones empresariales debido a su simplicidad de despliegue y su menor sobrecarga de comunicación interna \cite{aldebagy2018comparative}. Por su parte, la arquitectura modular --en particular bajo el paradigma de microservicios-- promueve la descomposición del sistema en componentes independientes, favoreciendo la escalabilidad y el mantenimiento, pero introduciendo mecanismos adicionales de comunicación, orquestación y coordinación entre servicios \cite{su2024microservice}. Esta tensión entre modularidad y eficiencia ha sido documentada en estudios comparativos que muestran resultados no siempre convergentes: mientras algunos trabajos reportan ahorros energéticos derivados de una mejor utilización de recursos en escenarios de alta demanda \cite{capuano2026comparative}, otros advierten que la sobrecarga de coordinación entre servicios puede incrementar el consumo total de energía, especialmente en sistemas de menor tamaño o fuertemente acoplados \cite{jagroep2017extending}.

Este proyecto se sitúa en el campo de la ingeniería de sistemas, la arquitectura de software y la eficiencia energética, un área de estudio que se ha consolidado bajo el concepto de \textit{Green Software Engineering} o ingeniería de software sostenible \cite{capra2012green}. La motivación de esta investigación surge de la necesidad de comprender, con evidencia empírica y no únicamente con supuestos de diseño, en qué medida la elección entre una arquitectura monolítica y una arquitectura modular incide sobre el consumo de recursos computacionales y, en consecuencia, sobre la huella energética del software producido.

El presente documento constituye la base metodológica de dicha investigación. En él se plantea la idea inicial de investigación, se delimita el problema, se identifican los síntomas, causas y consecuencias asociados, se construye una narrativa argumentativa del planteamiento del problema, se revisan los antecedentes académicos relevantes, se formula la pregunta de investigación, se establecen los objetivos y se plantean la hipótesis y las variables que orientarán las etapas posteriores del proyecto, entre ellas el diseño experimental.
